% R2 — Introduction: study question and contribution
We investigate whether a compact predictive skeleton model can produce useful representations from a small collection of gait videos. The model learns to predict hidden features of pose sequences without condition labels; a separate linear classifier then tests whether the frozen features contain information associated with the collection labels. Each encoder is trained on only 24 source recordings. Keeping all clips from a source together allows us to evaluate recordings that were excluded from both encoder and classifier fitting.

The main contribution is an empirical demonstration of useful trained features in this limited-data setting. S-JEPA exceeds a corrected Random Forest using handcrafted joint-angle summaries under both source-weighted and clip-weighted scoring. A mean-pose comparison helps assess how much can already be predicted from average pose without frame order, while a visibility-only comparison examines information available in the pose detector's outputs. Together, these experiments provide a grounded case study of representation learning for markerless behavioral analysis and identify the comparisons needed before interpreting the features as measures of neurological gait impairment.
